\begin{lemma}\label{lem:octic-315-cusp-image}
Let $F=\mathbb Q(\zeta_3)$, $X=X_1(3,15)$, and
$J=\operatorname{Jac}(X)$. The group $C\subseteq J(F)$
generated by $F$-rational cuspidal divisors has invariant factors
\[
 C\simeq[3,3,24,96,96],\qquad \#C=1990656,
\]
and $3J(F)\subseteq C$. Reduction identifies this same group
at residue fields $\mathbb F_4,\mathbb F_7,\mathbb F_{13}$,
and $\mathbb F_{19}$.
\end{lemma}
\begin{proof}
By \cite[Theorem~4.1]{DS}, $J$ has rank zero over $F$.
All cusps split over $L=\mathbb Q(\zeta_{15})$.
The prime $7$ splits in $F$, and its Frobenius generates
$\operatorname{Gal}(L/F)$, of order $4$.
Thus the cusp signature $1^{16},4^4$ at $7$ gives exactly
the global $F$-orbits. Their differences from an $F$-rational
cusp generate $C$.

For an explicit identification at the other primes, use the
normalized universal curve and sections
\[
 E:y^2+a_1xy+a_3y=x^3,\qquad P=(0,0),\qquad Q=(-t,t^2).
\]
The functions
\[
 a_3,\quad a_1^3-27a_3,\quad
 x(kQ+jP),\quad y(kQ+jP),\quad
 x(kQ)-x(kQ+P),
 \qquad 1\leq k\leq7,\quad 0\leq j\leq2,
\]
are integral units on the open modular curve away from $15$.
Indeed the discriminant is $a_3^3(a_1^3-27a_3)$; the
independence of $P$ and $5Q$ excludes all zeros and poles
of the displayed section coordinates and differences.
The normalization is integral: translate the flex $P$ to
the origin and its tangent to $y=0$; the value
$y(Q)/x(Q)^2$ is a unit and supplies the canonical scale.
The valuation vectors therefore commute with good reduction
and are constant on the global cusp orbits.

At $7$ these vectors distinguish all twenty global cusp
orbits. At the four residue fields they give the same
twenty blocks of degrees $1^{16},4^4$. At $\mathbb F_4$
and $\mathbb F_{19}$ each quartic block is the sum of two
quadratic closed cusps. Their exact class-group images
have the displayed invariant factors at all four fields.
The kernels of the common free cusp-divisor maps are
checked equal, so the identifications use the same global
divisors, independently of finite-group generator choices.

Good odd-prime unramified reduction is injective on torsion.
The exact Jacobian orders at $7,13,19$ are respectively
\[
149299200,\qquad12230590464,\qquad247669456896,
\]
whose gcd is $5971968=3\#C$. Hence $\#J(F)$ divides
$3\#C$. Since $C\subseteq J(F)$, the quotient $J(F)/C$
has order dividing $3$, proving $3J(F)\subseteq C$.
No assertion that $J(F)=C$ is needed.
\end{proof}
